\section{Introduction}
\label{sec:intro}

\subsection{What is actually here}

Strong-coupling expansions of Hamiltonian lattice gauge theory are old.
What is not old is the object this program organizes them around.

Write the cellular chain complex of the cubic lattice as
$C_3 \xrightarrow{\partial_3} C_2 \xrightarrow{\partial_2} C_1$, with
$C_2$ spanned by plaquettes. The subspace
\begin{equation}
  Z_2 = \ker \partial_2 \subset C_2
\end{equation}
--- the closed plaquette chains --- is the carrier of the strong-coupling
band. It is a homological object, fixed before any perturbation theory
is done, and the entire content of Parts~\ref{part:exact} and
\ref{part:fixed} is a description of what the Kogut--Susskind Hamiltonian
does to it, order by order and rank by rank.

Three things follow from taking that object seriously, and they are the
three claims of this document.

\emph{First: the coefficients are exact rational functions of the rank.}
Not asymptotic in $N$, not fitted, not computed rank by rank and
extrapolated. The second-order hopping amplitude is
\begin{equation}
  t_N = \frac{2N\,(N^2-4)}{(N^2-1)(2N^2-1)(4N^2-9)},
  \qquad t_2 = 0, \quad t_3 = \tfrac{5}{612},
  \label{eq:intro-tN}
\end{equation}
and the four channel weights behind it are order-two Weingarten values,
so \eqref{eq:intro-tN} imports nothing from any expansion. The factor
$N^2-4$ is not a coincidence of small rank: it is why $\SU(2)$ has no
band at this order at all. The same discipline runs through fourth
order, through the rank grading of \S\ref{sec:grading}, and into a
symbolic engine that works over $\Q(N)$ with $N$ an indeterminate rather
than a number.

\emph{Second: the rank grading has exact structure, and it comes from
cancellation.} For the $C$-parity splitting of the one-plaquette
diagonal, every even order carries the same single power of $N$ once the
coupling is written in the planar variable, and the leading coefficients
of the four contributing intermediate channels are in the ratio
$(-4, +2, +1, +1)$ and \emph{sum to zero}. The grading is a cancellation
theorem, not a power count. Complementarily, the odd orders vanish
identically for even $N$, and for odd $N$ the first surviving one is
order $N-2$, with a closed-form vertex. Neither statement is visible from
any single rank.

\emph{Third: the construction at fixed lattice spacing is complete, and
it does not reach the continuum.} \S\ref{sec:wilson} states an actual
infinite-volume Wilson transfer limit with a complete physical odd band
and a literal plaquette-source frame that is onto it. \S\ref{sec:breaks}
states why that is not a continuum mass gap, in six independent ways,
five of them with finite counterexamples. The most important is
structural rather than technical: the theorem holds on
$\abs{u} \le u_*/(10\,022\,400\,000 N)$ and the continuum trajectory runs
to $u \to \infty$, so the two are disjoint and no sharpening of the
constant changes that.

\subsection{Why the document is shaped this way}

Alex, the maintainer of this program, asked for a paper that would serve
as the base for further work on the mass gap rather than as a report of
work completed. That request has consequences for the form.

It means that failure is content. \S\ref{sec:breaks} and
\S\ref{sec:register} exist because a route this program has already
closed is more valuable to the next attempt than a route it has already
opened, and because at least four arguments were written down here as
complete and then found not to be. It means the route register in
\S\ref{sec:obligations} reproduces \RouteDead{} dead routes at length,
including their measured costs, rather than listing only the \RouteLive{}
live ones. And it means every count in the prose is a macro generated
from the repository's ledgers, so that a stale number is a build failure
rather than a quiet misstatement.

It also means the document distinguishes, everywhere and visibly, three
things that are routinely conflated: what a statement mathematically
establishes, what artefact backs it, and what a machine checks. A result
can be fully proved analytically and carry the weakest machine badge.
A Lean theorem can be checked by the kernel and prove something
trivial. This program has made both errors and the notation of
\S\ref{sec:howto} is the correction.

\subsection{Relation to the literature}

The external comparisons are recorded as an evidence map rather than a
citation list: \LitPapers{} indexed papers, \LitEdges{} recorded
relations, of which \LitVerified{} are verified against a named claim
here. A published paper is not authority in this program --- it is
unchecked until something checks it, exactly like an internal document.
What makes an external agreement evidence is that it was produced
without knowledge of this program, and \S\ref{sec:external} reports the
comparisons that hold and, separately, the two that contradict.

\subsection{What is claimed, and what is not}

Claimed: the contents of Part~\ref{part:exact}, unconditionally, at the
tiers recorded; and the contents of Part~\ref{part:fixed} on the stated
interval and under the stated hypotheses.

Not claimed: the existence of continuum four-dimensional Yang--Mills
theory; a continuum mass gap; Poincar\'e covariance; restoration of
rotational invariance from the hypercubic lattice; or the construction of
a continuum measure. None of these has an argument in this document, and
\S\ref{sec:breaks} explains for each of the six obstructions why the
present machinery does not supply one.
